\chapter{Construção das Variáveis}
\label{app:variaveis}

Este apêndice reúne as definições das variáveis usadas nos capítulos e indica
a coluna correspondente nos arquivos de dados (Apêndice~\ref{app:pastas}). As
séries são diárias, de 24/03/2021 a 03/03/2026. O preço de fechamento do
Bitcoin (Binance, par BTCUSDT) e o DVOL (Deribit) da data $t$ são os da vela
diária que abre às 00:00 UTC de $t$; os dois são observados no fim do dia $t$.

\section{Retornos}

O log-retorno diário,
\[
r_t = \ln\left(\frac{P_t}{P_{t-1}}\right),
\]
em que $P_t$ é o preço de fechamento, entra no cálculo da volatilidade
histórica e, como variável explicativa, nos log-retornos de 1, 5 e 30 dias
(colunas \texttt{ret} ou \texttt{ret\_1d}, \texttt{ret\_acum\_5d} e
\texttt{ret\_acum\_30d}). O retorno futuro é o retorno simples
acumulado em $h$ dias,
\[
R_{t+h} = \frac{P_{t+h}}{P_t} - 1, \qquad h \in \{1, 5, 10, 20, 30, 60\},
\]
nas colunas \texttt{ret\_fut\_1d} a \texttt{ret\_fut\_60d}. Ele é a variável
dependente dos Capítulos~\ref{chap:retornos_linear}
e~\ref{chap:retornos_nao_linear} e nunca entra como variável explicativa. Os
dois tipos de retorno não se misturam: a volatilidade usa log-retornos, e os
retornos futuros e as estratégias do Apêndice~\ref{app:estrategias}, retornos
simples.

\section{Volatilidade histórica}

A volatilidade histórica de $k$ dias, em \% a.a., é
\[
\text{VH}^{(k)}_t = 100 \sqrt{\frac{365}{k} \sum_{j=0}^{k-1} r_{t-j}^2},
\]
sem subtrair a média dos retornos, com anualização por 365 dias porque o
mercado funciona todos os dias. Com dados diários, a VH é o estimador da
volatilidade realizada (RV) que entra na definição do prêmio
(Seção~\ref{sec:refer-vrp-teorico}). A janela de 30 dias aparece em duas
posições:
\begin{itemize}
    \item retrospectiva, $\text{VH}_{t-29:t}$, com os retornos de $t-29$ a
          $t$, conhecida em $t$ (colunas \texttt{rv\_30d} e \texttt{vh\_30d});
    \item prospectiva, $\text{VH}_{t+1:t+30}$, com os retornos de $t+1$ a
          $t+30$, conhecida só em $t+30$ (coluna \texttt{rv\_30d\_fut}). Por
          construção, a VH prospectiva de $t$ é a retrospectiva de $t+30$.
\end{itemize}
As janelas de 1, 5, 60 e 90 dias entram só como variáveis explicativas
(\texttt{vh\_1d}, \texttt{vh\_5d}, \texttt{vh\_60d} e \texttt{vh\_90d}).

\section{Volatilidade implícita}

A volatilidade implícita de 30 dias, $\text{IV}_t$, é o fechamento diário do
índice DVOL da Deribit, em \% a.a. (coluna \texttt{iv\_30d}). O DVOL já tem
horizonte de 30 dias (Seção~\ref{sec:refer-dvol}), e a dissertação não
interpola vencimentos.

\section{Prêmio de risco de volatilidade}

A Tabela~\ref{tab:var-premio} resume as três medidas do prêmio. Com a
convenção $\text{VH} - \text{IV}$, um valor negativo indica que a IV supera a
volatilidade histórica, e um prêmio maior corresponde a um valor mais
negativo.

\begin{table}[H]
\centering
\caption{Medidas do prêmio de risco de volatilidade: definição, coluna nos
arquivos de dados e data em que a medida é conhecida.}
\label{tab:var-premio}
\small
\begin{tabular}{llll}
\toprule
Símbolo & Definição & Coluna & Conhecido em \\
\midrule
$\text{BVRP}_t$ & $\text{VH}_{t+1:t+30} - \text{IV}_t$ & \texttt{bvrp\_30d\_fut} & $t+30$ \\
$\text{BVRP}^{\text{proxy}}_t$ & $\text{VH}_{t-29:t} - \text{IV}_t$ & \texttt{vrp\_30d} & $t$ \\
$\text{BVRP}_{t-30}$ & $\text{VH}_{t-29:t} - \text{IV}_{t-30}$ & \texttt{bvrp\_realizado\_defasado} & $t$ \\
\bottomrule
\end{tabular}
\end{table}

O BVRP é a variável-alvo do Capítulo~\ref{chap:previsao_bvrp} (no arquivo da
amostra de modelagem, coluna \texttt{alvo\_bvrp\_30d\_fut}). A proxy supõe um
passeio aleatório sem deriva para a volatilidade histórica e entra como
variável explicativa. O prêmio realizado defasado é o BVRP de $t-30$, que se
realiza em $t$; ele é igual à proxy mais a variação da IV em 30 dias,
$\text{IV}_t - \text{IV}_{t-30}$.

\section{Variáveis explicativas}

As 17 variáveis explicativas do Capítulo~\ref{chap:previsao_bvrp} estão no
dicionário da Tabela~\ref{tab:prev-dicionario}, com fórmula, janela e
tratamento. Todas são funções dos preços e do DVOL até $t$, calculadas por
uma única função do script \texttt{code/build\_ml\_dataset\_T4.py}. O teste
\texttt{code/test\_T4\_vazamento.py} corta os dados brutos em datas sorteadas,
recalcula as variáveis e confere que nenhuma usa informação posterior a $t$.

\section{Regimes de volatilidade}

O regime de alta volatilidade é
\[
D_t = \mathbf{1}\left\{\text{VH}_{t-29:t} > c\right\}.
\]
No corte principal, $c = 37{,}3\%$ a.a. é o ponto em que a probabilidade a
posteriori do componente de alta volatilidade passa de 1/2 numa mistura de
duas normais ajustada ao logaritmo da VH de 30 dias, estimada só na primeira
janela de estimação, de 23/04/2021 a 22/04/2023
(Seção~\ref{sec:dados-regimes}). Como robustez, a dissertação reestima o corte
pelo mesmo método em cada origem de previsão, só com os dados até ela
(Seção~\ref{sec:met-regimes}); as datas anteriores à primeira origem usam o
corte fixo. O arquivo \texttt{data/regimes\_T9.csv} traz,
para cada data, os dois cortes (\texttt{corte\_fixo} e
\texttt{corte\_expansivo\_vigente}) e as duas classificações
(\texttt{regime\_alta\_fixo} e \texttt{regime\_alta\_expansivo}).

\section{Informação disponível em \texorpdfstring{$t$}{t}}

As variáveis explicativas de cada data são as da própria data, todas
conhecidas no fim de $t$. O alvo de uma data $s$ olha $h$ dias à frente e só
é conhecido em $s+h$; por isso, o modelo estimado na origem $t$ usa apenas as
datas $s \le t - h$, e a validação cruzada repete o mesmo embargo dentro do
treino (Seção~\ref{sec:met-amostras}). O BVRP nunca entra como variável
explicativa.
